\section{Between-Subjects Survey Experiment}
\label{sec:method-survey}
We designed a between-subjects survey experiment to investigate how exposure to minority views affects readers' moral judgments, reasoning, and consideration of alternative perspectives (\textbf{RQ4}).

\subsection{Survey Design}
\label{subsec:survey-design}

\paragraph{Design and conditions}
We conducted a between-subjects design crossing four moral-dilemma scenarios with two comment-display conditions.
Participants were randomly assigned to one of two comment-display conditions:
(1) a majority--majority condition, in which participants saw the two highest-scoring comments supporting the same side, or
(2) a majority--minority condition, in which participants saw the highest-scoring comment supporting each side.

\paragraph{Scenarios and stimuli}
We selected four posts from the AITA community that showed high levels of agreement among commenters: two with a majority \textit{YA} judgment and two with a majority \textit{NA} judgment. The four posts covered distinct topics: financial responsibility between friends, family and parenting responsibilities, punctuality and obligations in friendships, and professional conduct in the workplace.

% TODO: Describe thread selection and the presentation of replies, scores,
% and alignment information.

\paragraph{Participants and procedure}
% TODO: Report recruitment, eligibility, and exclusion criteria.
We have recruited 38 participants to date: 18 in the majority--majority condition and 20 in the majority--minority condition.

Participants were randomly assigned to one of four scenarios and one of two comment-display conditions. The procedure consisted of three phases. First, participants read an AITA post and its discussion thread. Second, they reported their moral judgment and briefly explained their reasoning. Finally, they completed self-report measures of comment influence, agreement, reflection, perspective broadening, and comment-display preference.

\paragraph{Measures}

The primary outcome was participants’ moral judgment of the original poster, measured on an eight-point scale from 1 (Definitely YA) to 8 (Definitely NA). Participants also provided an open-ended explanation of at least two sentences. For each featured comment, participants rated perceived influence (1–5), agreement (1–5), and perspective elicitation (1–5). They also indicated whether Comment 1, Comment 2, both equally, or neither contributed more to their judgment. Discussion-level perspective broadening was measured using three seven-point items assessing whether the discussion encouraged consideration of both sides, introduced new considerations, and prompted engagement with reasons supporting an alternative judgment. Finally, participants indicated their preferred comment-display format: two highly upvoted comments supporting the same judgment, two highly upvoted comments representing opposing judgments, or neither.

% TODO: Define judgment certainty and document its scale or derivation.
% TODO: Confirm that Comment 1 is the shared majority comment and Comment 2
% is the second majority comment or the minority comment across conditions.
% TODO: Add a survey-procedure diagram to the Appendix.

\subsection{Survey Results}
\label{subsec:survey-results}

The current sample comprises 38 participants. Table~\ref{tab:survey-results} summarizes the preliminary condition means. These results address whether exposure to a minority thread broadens readers' perspectives and informs their judgments (\textbf{RQ4}).

% TODO: Specify statistical tests, test statistics, exact p-values, effect
% sizes, uncertainty intervals, and how scenario differences are handled.
% Significance statements below reflect the supplied preliminary analysis.

\paragraph{Perspective broadening and consideration}
Overall perspective broadening did not differ significantly between the majority--majority ($M=4.43$) and majority--minority ($M=4.55$) conditions. Perspective consideration elicited by Comment 2 was descriptively higher in the majority--minority condition ($M=2.75$ versus $M=2.11$), but this difference was not statistically significant. Thus, the current results do not establish greater perspective broadening from exposure to a minority thread.

\paragraph{Perceived influence and agreement}
Comment 2 received significantly lower agreement in the majority--minority condition ($M=2.75$) than in the majority--majority condition ($M=4.22$). Its perceived influence was also descriptively lower ($M=2.20$ versus $M=2.28$). For Comment 1, mean influence was $2.90$ versus $2.56$, and mean agreement was $4.00$ versus $4.28$, respectively. These descriptive differences should not be interpreted as evidence of condition effects without the corresponding statistical tests.

\paragraph{Judgment certainty}
Judgment certainty was significantly lower in the majority--minority condition ($M=2.50$) than in the majority--majority condition ($M=3.61$). Taken together, the preliminary results indicate lower certainty and lower agreement with Comment 2 when a minority thread is included, while providing no statistically significant evidence of increased perspective broadening.

\begin{table}[t]
  \centering
  \small
  \caption{Preliminary survey means for the majority--majority (MM; $n=18$) and majority--minority (Mm; $n=20$) conditions. Differences are Mm minus MM. Statistical significance is described in the text; test details remain to be added.}
  \label{tab:survey-results}
  \begin{tabular}{@{}lrrr@{}}
    \toprule
    Measure & MM & Mm & $\Delta$ \\
    \midrule
    Judgment certainty & 3.61 & 2.50 & $-1.11$ \\
    Comment 1 influence & 2.56 & 2.90 & $+0.34$ \\
    Comment 2 influence & 2.28 & 2.20 & $-0.08$ \\
    Comment 1 agreement & 4.28 & 4.00 & $-0.28$ \\
    Comment 2 agreement & 4.22 & 2.75 & $-1.47$ \\
    Comment 1 perspective consideration & 2.06 & 2.25 & $+0.19$ \\
    Comment 2 perspective consideration & 2.11 & 2.75 & $+0.64$ \\
    Overall perspective broadening & 4.43 & 4.55 & $+0.12$ \\
    \bottomrule
  \end{tabular}
\end{table}
